% ============================================================
%  Глава 1 — 2D Velocity Extended Kalman Filter
%  Файл: pi_nodes/filters/velocity_ekf.py
% ============================================================
\chapter{2D Velocity Extended Kalman Filter}
\label{ch:velocity_ekf}

\begin{filebox}[title={Файл исходного кода}]
\filepath{pi\_nodes/filters/velocity\_ekf.py} \quad (251 строка, Python)
\end{filebox}

\section{Постановка задачи}

Для точного позиционирования гусеничного робота необходимо объединить два
источника скорости:

\begin{itemize}
  \item \textbf{Колёсная одометрия} --- стабильна, не имеет дрейфа, но не учитывает
    проскальзывание гусениц на неровном покрытии.
  \item \textbf{Акселерометр (однократное интегрирование)} --- улавливает реальное
    скольжение, но накапливает ошибку интегрирования (дрейф).
\end{itemize}

Оба источника объединяются \textbf{фильтром Калмана}, который автоматически
взвешивает вклад каждого по его статистическим шумовым характеристикам.

\section{Вектор состояния и модель}

\subsection{Вектор состояния}

Состояние фильтра --- скорость робота в \emph{мировой} системе координат:

\begin{equation}
  \vect{x} = \begin{pmatrix} v_x \\ v_y \end{pmatrix} \in \mathbb{R}^2
  \label{eq:ekf_state}
\end{equation}

Позиция получается последовательным интегрированием фильтрованной скорости.

\subsection{Матрица ковариации}

Поскольку numpy не используется, матрица $2\times 2$ хранится в четырёх скалярах:

\begin{equation}
  \mat{P} = \begin{pmatrix} p_{00} & p_{01} \\ p_{10} & p_{11} \end{pmatrix}
  \label{eq:cov_matrix}
\end{equation}

Начальное значение: $\mat{P} = 0.01 \cdot \mat{I}$.

\section{Шаг предсказания (Prediction)}
\label{sec:ekf_predict}

\begin{filebox}[title={\filepath{velocity\_ekf.py}, строки 90--120 --- predict\_wheel()}]
\begin{lstlisting}[style=pythonstyle, numbers=left, firstnumber=104]
# Wheel-predicted velocity in world frame
vx_wheel = vx_cmd * cos(theta)
vy_wheel = vx_cmd * sin(theta)

# Prediction: state transitions toward wheel command
alpha = 0.5
self.vx = (1 - alpha) * self.vx + alpha * vx_wheel
self.vy = (1 - alpha) * self.vy + alpha * vy_wheel

# Process noise: covariance grows
self._p00 += self.q_wheel * dt
self._p11 += self.q_wheel * dt

# Integrate position using current fused velocity
self.x += self.vx * dt
self.y += self.vy * dt
\end{lstlisting}
\end{filebox}

\textbf{Математика предсказания:}

Команда колёс задаёт вперёд\-направленную скорость $v_{\text{cmd}}$ в теле робота.
Поворот на угол курса $\theta$ (из IMU) переводит её в мировую систему координат:

\begin{equation}
  \begin{pmatrix} v_x^{\text{wheel}} \\ v_y^{\text{wheel}} \end{pmatrix}
  = v_{\text{cmd}} \begin{pmatrix} \cos\theta \\ \sin\theta \end{pmatrix}
  \label{eq:wheel_world}
\end{equation}

Смешивание с текущим состоянием ($\alpha = 0.5$):

\begin{equation}
  \vect{x}^- = (1-\alpha)\,\vect{x} + \alpha\,\vect{x}^{\text{wheel}}
  \label{eq:blend}
\end{equation}

Рост ковариации (процессный шум $q_{\text{wheel}} = 0.01$):

\begin{equation}
  p_{00}^- = p_{00} + q_{\text{wheel}} \cdot \Delta t, \quad
  p_{11}^- = p_{11} + q_{\text{wheel}} \cdot \Delta t
  \label{eq:cov_grow}
\end{equation}

Интегрирование позиции (правило Эйлера):

\begin{equation}
  \boxed{x \mathrel{+}= v_x \cdot \Delta t, \quad y \mathrel{+}= v_y \cdot \Delta t}
  \label{eq:pos_integrate}
\end{equation}

\section{Шаг коррекции (Update) по акселерометру}
\label{sec:ekf_update}

\begin{filebox}[title={\filepath{velocity\_ekf.py}, строки 122--173 --- update\_accel()}]
\begin{lstlisting}[style=pythonstyle, numbers=left, firstnumber=135]
# Integrate accel to get velocity measurement
self._accel_vx += la_wx * dt
self._accel_vy += la_wy * dt

# Decay accel velocity to prevent unbounded drift
self._accel_vx *= 0.98
self._accel_vy *= 0.98

# Innovation: y = z - H*x  (H = I)
yx = self._accel_vx - self.vx
yy = self._accel_vy - self.vy

# S = P + R (innovation covariance)
r = self.r_accel + self.q_accel * dt
s00 = self._p00 + r
s11 = self._p11 + r

# Kalman gain: K = P * S^-1
k00 = self._p00 / s00
k11 = self._p11 / s11

# State update
self.vx += k00 * yx
self.vy += k11 * yy

# Covariance update: P = (I - K*H) * P
self._p00 *= (1 - k00)
self._p11 *= (1 - k11)
self._p01 *= (1 - 0.5 * (k00 + k11))
self._p10 *= (1 - 0.5 * (k00 + k11))
\end{lstlisting}
\end{filebox}

\textbf{Математика коррекции:}

Однократное интегрирование ускорения (измерение скорости):

\begin{equation}
  \vect{v}^{\text{accel}}_{k} = \vect{v}^{\text{accel}}_{k-1}
  + \vect{a}^{\text{world}} \cdot \Delta t
  \label{eq:accel_integrate}
\end{equation}

Экспоненциальное затухание для борьбы с дрейфом:

\begin{equation}
  \vect{v}^{\text{accel}} \leftarrow 0.98 \cdot \vect{v}^{\text{accel}}
  \label{eq:decay}
\end{equation}

Инновация (невязка):

\begin{equation}
  \vect{y} = \vect{z} - \mat{H}\,\vect{x}^-, \quad \mat{H} = \mat{I}
  \quad\Rightarrow\quad
  \vect{y} = \vect{v}^{\text{accel}} - \vect{v}^{\text{state}}
  \label{eq:innovation}
\end{equation}

Ковариация инновации:

\begin{equation}
  \mat{S} = \mat{P}^- + \mat{R}, \quad
  R = r_{\text{accel}} + q_{\text{accel}} \cdot \Delta t
  \label{eq:S_matrix}
\end{equation}

\textbf{Усиление Калмана} (для диагональной $\mat{S}$, т.е. поэлементно):

\begin{equation}
  \boxed{K_{ii} = \frac{p_{ii}^-}{S_{ii}}, \quad i \in \{0,1\}}
  \label{eq:kalman_gain}
\end{equation}

\textbf{Обновление состояния:}

\begin{equation}
  \vect{x} = \vect{x}^- + \mat{K}\,\vect{y}
  \label{eq:state_update}
\end{equation}

\textbf{Обновление ковариации} (формула Джозефа в диагональном виде):

\begin{equation}
  \boxed{p_{ii} \leftarrow (1 - K_{ii})\,p_{ii}}
  \label{eq:cov_update}
\end{equation}

\section{ZUPT --- обновление нулевой скорости}

\begin{filebox}[title={\filepath{velocity\_ekf.py}, строки 181--207 --- update\_zupt()}]
\begin{lstlisting}[style=pythonstyle, numbers=left, firstnumber=186]
# Kalman gain (measurement = zero, R = r_zupt)
s00 = self._p00 + r
s11 = self._p11 + r
k00 = self._p00 / s00
k11 = self._p11 / s11

# Update: z = [0, 0], innovation = -[vx, vy]
self.vx -= k00 * self.vx
self.vy -= k11 * self.vy

# Covariance update
self._p00 *= (1 - k00)
self._p11 *= (1 - k11)
\end{lstlisting}
\end{filebox}

Когда робот стоит, измерение $\vect{z} = \vect{0}$ с очень малым шумом
$r_{\text{zupt}} = 10^{-4}$. Инновация: $\vect{y} = \vect{0} - \vect{v} = -\vect{v}$.
Это приводит к быстрому обнулению скорости через стандартный шаг Калмана.

\section{Неголономное ограничение}
\label{sec:nonholomic}

\begin{filebox}[title={\filepath{velocity\_ekf.py}, строки 209--228 --- apply\_nonholonomic()}]
\begin{lstlisting}[style=pythonstyle, numbers=left, firstnumber=213]
cy, sy = cos(theta), sin(theta)
fwd = self.vx * cy + self.vy * sy    # forward component
lat = -self.vx * sy + self.vy * cy   # lateral component

lat *= self._lateral_decay           # kill 70% of lateral vel

self.vx = fwd * cy - lat * sy
self.vy = fwd * sy + lat * cy
\end{lstlisting}
\end{filebox}

Гусеничный робот \textbf{не может двигаться боком}. После шага Калмана
производится проекция скорости на направление движения:

\begin{align}
  v_{\text{fwd}} &= v_x \cos\theta + v_y \sin\theta \label{eq:fwd} \\
  v_{\text{lat}} &= -v_x \sin\theta + v_y \cos\theta \label{eq:lat}
\end{align}

Боковая составляющая подавляется: $v_{\text{lat}} \leftarrow 0.3 \cdot v_{\text{lat}}$.

Обратное преобразование в мировую СК (поворот матрицей $\mat{R}^{-1} = \mat{R}\T$):

\begin{equation}
  \begin{pmatrix} v_x \\ v_y \end{pmatrix}
  = \begin{pmatrix} \cos\theta & -\sin\theta \\ \sin\theta & \cos\theta \end{pmatrix}
  \begin{pmatrix} v_{\text{fwd}} \\ v_{\text{lat}} \end{pmatrix}
  \label{eq:rot_back}
\end{equation}

\section{Модуль скорости}

\begin{filebox}[title={\filepath{velocity\_ekf.py}, строки 243--245}]
\begin{lstlisting}[style=pythonstyle, numbers=left, firstnumber=243]
@property
def velocity_magnitude(self) -> float:
    return sqrt(self.vx * self.vx + self.vy * self.vy)
\end{lstlisting}
\end{filebox}

\begin{equation}
  |\vect{v}| = \sqrt{v_x^2 + v_y^2}
  \label{eq:vmag}
\end{equation}

\section{Теоретическое обоснование: связь с ТАУ}

\subsection{Модель в пространстве состояний}

Velocity EKF реализует дискретную модель в пространстве состояний
(см.\ гл.~\ref{ch:theory}, формула~\eqref{eq:state_space}):

\begin{equation}
  \underbrace{\begin{pmatrix} v_x \\ v_y \end{pmatrix}_{k}}_{\vect{x}_k}
  = \underbrace{\begin{pmatrix} 1-\alpha & 0 \\ 0 & 1-\alpha \end{pmatrix}}_{\mat{F}}
  \underbrace{\begin{pmatrix} v_x \\ v_y \end{pmatrix}_{k-1}}_{\vect{x}_{k-1}}
  + \underbrace{\begin{pmatrix} \alpha\cos\theta & 0 \\ 0 & \alpha\sin\theta \end{pmatrix}}_{\mat{B}}
  \underbrace{\begin{pmatrix} v_{\text{cmd}} \\ v_{\text{cmd}} \end{pmatrix}}_{\vect{u}}
  \label{eq:vekf_ss}
\end{equation}

\subsection{Устойчивость фильтра}

Собственные числа матрицы перехода $\mat{F}$:
$\lambda_{1,2} = 1 - \alpha = 0.5$.
Поскольку $|\lambda_i| < 1$ для дискретной системы,
фильтр \textbf{асимптотически устойчив} (ошибка сходится к нулю).

Это дискретный аналог критерия устойчивости~\eqref{eq:stability_eigenvalues}:
для дискретных систем все собственные числа должны лежать
внутри единичного круга $|\lambda_i| < 1$.

\subsection{Оптимальность усиления Калмана}

Усиление Калмана $K_{ii} = p_{ii}/(p_{ii} + R)$
минимизирует апостериорную дисперсию оценки:

\begin{equation}
  \frac{\partial}{\partial K}\,E\!\left[(x - \hat{x})^2\right] = 0
  \quad\Rightarrow\quad K^* = \frac{P^-}{P^- + R}
  \label{eq:optimal_K}
\end{equation}

При $R \to 0$ (идеальный датчик): $K \to 1$, фильтр доверяет измерению.
При $R \to \infty$ (шумный датчик): $K \to 0$, фильтр доверяет модели.

\section{Сводная таблица параметров}

\begin{center}
\begin{tabular}{llll}
\toprule
\textbf{Параметр} & \textbf{Символ} & \textbf{Значение} & \textbf{Физический смысл}\\
\midrule
\texttt{q\_wheel} & $q_w$ & $0.01\,(m/s)^2$ & Шум процесса (колёса)\\
\texttt{q\_accel} & $q_a$ & $0.05\,(m/s)^2$ & Скорость дрейфа акселерометра\\
\texttt{r\_accel} & $r_a$ & $0.1\,(m/s)^2$ & Шум измерения акселерометра\\
\texttt{r\_zupt}  & $r_z$ & $10^{-4}\,(m/s)^2$ & Доверие к нулевой скорости\\
\texttt{lateral\_decay} & --- & $0.3$ & Подавление бокового движения\\
\bottomrule
\end{tabular}
\end{center}
